% docs/paper.tex
% PACM v1.0.0 — Publication-Ready LaTeX Manuscript
% Session 11: CoT trace extraction, scope-flag friction boost, real-model sweep
%
% Compile with:
%   pdflatex paper.tex
%   bibtex paper
%   pdflatex paper.tex
%   pdflatex paper.tex
%
% Or with latexmk:
%   latexmk -pdf paper.tex

\documentclass[11pt,letterpaper]{article}

% ── Packages ──────────────────────────────────────────────────────────────────
\usepackage[utf8]{inputenc}
\usepackage[T1]{fontenc}
\usepackage{lmodern}
\usepackage[margin=1in]{geometry}
\usepackage{hyperref}
\usepackage{booktabs}
\usepackage{tabularx}
\usepackage{multirow}
\usepackage{amsmath}
\usepackage{amssymb}
\usepackage{algorithm}
\usepackage{algpseudocode}
\usepackage{graphicx}
\usepackage{listings}
\usepackage{xcolor}
\usepackage{enumitem}
\usepackage{microtype}
\usepackage{threeparttable}
\usepackage{caption}
\usepackage{subcaption}
\usepackage{natbib}
\usepackage{times}

% ── Hyperref configuration ────────────────────────────────────────────────────
\hypersetup{
  colorlinks=true,
  linkcolor=blue!70!black,
  citecolor=green!50!black,
  urlcolor=blue!60!black,
  pdfauthor={PACM Research Team},
  pdftitle={Plan-Action Consistency Monitoring (PACM)},
  pdfkeywords={prompt injection, LLM agents, MCP, runtime security,
               plan-action consistency, divergence scoring}
}

% ── Code listing style ────────────────────────────────────────────────────────
\lstset{
  basicstyle=\ttfamily\small,
  keywordstyle=\color{blue}\bfseries,
  commentstyle=\color{gray}\itshape,
  stringstyle=\color{green!60!black},
  backgroundcolor=\color{gray!5},
  frame=single,
  framesep=4pt,
  rulecolor=\color{gray!40},
  breaklines=true,
  captionpos=b,
  numbers=left,
  numberstyle=\tiny\color{gray},
  numbersep=5pt
}

% ── Custom commands ───────────────────────────────────────────────────────────
\newcommand{\pacm}{\textsc{PACM}}
\newcommand{\mcp}{\textsc{MCP}}
\newcommand{\abr}{\textit{ABR}}
\newcommand{\fpr}{\textit{FPR}}
\newcommand{\btsr}{\textit{BTSR}}
\newcommand{\calD}{\mathcal{D}}
\newcommand{\calF}{\mathcal{F}}
\newcommand{\todo}[1]{\textcolor{red}{\textbf{TODO:} #1}}

% ── Title & author ────────────────────────────────────────────────────────────
\title{%
  \textbf{Plan--Action Consistency Monitoring (\pacm{}):\\
  A Lightweight Runtime Defense Against Context-Aware\\
  Prompt Injection in Tool-Using LLM Agents}%
}

\author{%
  \textbf{PACM Research Team}\\
  \small{\texttt{pacm-research@example.com}}
}

\date{August 2026 \quad --- \quad v1.0.0 (arXiv preprint)}

% ── Document ──────────────────────────────────────────────────────────────────
\begin{document}

\maketitle

% ─────────────────────────────────────────────────────────────────────────────
\begin{abstract}
Large language model (LLM) agents that invoke external tools via the Model
Context Protocol (\mcp{}) are increasingly vulnerable to \emph{context-aware
prompt injection}, where adversarial content embedded in tool outputs hijacks
the agent's subsequent actions. We present \textbf{Plan--Action Consistency
Monitoring (\pacm{})}, a lightweight, protocol-agnostic runtime defense that
intercepts tool calls before execution and measures the \emph{divergence}
between the agent's declared reasoning intent and the proposed action
parameters. \pacm{} combines an embedding-based semantic scorer (TF-IDF or
\texttt{sentence-transformers}/all-MiniLM-L6-v2) with rule-based argument
scope checks and \mcp{} tool trust scores into a \textbf{trust-weighted
friction policy} $F = D \times (1.5 - 0.5T)$ that issues
\textsc{Allow}/\textsc{Flag}/\textsc{Block} decisions at runtime, without
requiring LLM calls or ground-truth labels at inference time. Evaluated on a
16-configuration benchmark across two agent harnesses (LangGraph and a
framework-free custom loop), two attack types, and three policy tiers,
\pacm{} achieves up to \textbf{54\% Attack Block Rate (ABR)} with
\textbf{0\% False Positive Rate (FPR)} on synthetic adversarial traces.
Alpha sensitivity analysis identifies $\alpha = 0.9$ as the optimal semantic
weight (separation = 0.479). Adaptive attacker stress tests (RQ5) reveal
consistent \textsc{Flag} escalation even for fully adaptive adversaries,
preserving a human oversight channel. Ablations confirm that trust modulation
and the sentence-transformers backend each contribute meaningfully to
detection. \pacm{} is implemented as a modular pre-tool hook requiring zero
changes to agent harness internals. Code, data, and reproduction scripts are
released at this repository.
\end{abstract}

% ─────────────────────────────────────────────────────────────────────────────
\section{Introduction}
\label{sec:intro}

The proliferation of LLM agents capable of executing multi-step tool-calling
workflows has introduced a new attack surface: \emph{prompt injection via tool
outputs}. When an agent reads from an external source (email inbox, web page,
database record) and that content contains adversarial instructions, the agent
may be induced to exfiltrate data, impersonate users, or execute unintended
system actions --- a class of attacks collectively termed
\emph{context-aware prompt injection} (CAPI)
\citep{Greshake2023NotWhatYouSignedUp}.

Existing defenses fall into three broad categories: (1)~\emph{input
sanitization} \citep{Liu2023PromptInjectionAttack}, (2)~\emph{structural
isolation} \citep{Debenedetti2025CaMeL}, and (3)~\emph{training-time
instruction hierarchies} \citep{Wallace2024InstructionHierarchy}. All three
approaches share a fundamental limitation: they are benchmarked on
context-independent tasks where the correct action is specified up front.
The 2026 survey by \citet{Wang2026AgentPI} explicitly identifies
context-dependent evaluation and the utility--security tradeoff as the
field's biggest open problems.

We present \pacm{}, which belongs to a fourth category --- \emph{runtime
behavioral monitoring} --- and extends it with a principled scoring framework
grounded in the \emph{consistency} between declared reasoning and actual
actions. Rather than classifying the content of tool outputs (brittle and
gameable), \pacm{} asks: \emph{does the tool call the agent is about to make
match what the agent just said it intended to do?} Injection-induced hijacks
manifest precisely as divergences between these two signals.

\paragraph{Contributions.}
\begin{enumerate}[leftmargin=*]
  \item \textbf{Divergence formalization:} We formalize the plan-action
    divergence problem and show that injection attacks reliably manifest as
    measurable divergences in both semantic embedding space and argument
    scope space (Section~\ref{sec:scoring}).

  \item \textbf{DivergenceScorer:} A dual-signal scorer combining semantic
    embedding similarity (TF-IDF or sentence-transformers) with argument
    scope rule checks, producing $D \in [0,1]$ with configurable weighting
    $\alpha$ (Section~\ref{sec:scoring}).

  \item \textbf{Trust-Weighted FrictionPolicy:} $F = \mathrm{clip}(D \times
    (1.5 - 0.5T),\, 0,\, 1)$, mapping friction to
    \textsc{Allow}/\textsc{Flag}/\textsc{Block} using configurable
    thresholds and \mcp{} tool trust scores $T$ (Section~\ref{sec:policy}).

  \item \textbf{RQ5 adaptive stress test:} We implement and evaluate five
    adaptive attack variants against \pacm{}, characterizing the monitor's
    own attack surface (Section~\ref{sec:adaptive}).

  \item \textbf{Open reproducibility:} A one-click \texttt{reproduce\_all.py}
    script replicates all 8 experimental steps; 245 unit tests verify all
    components (Appendix~\ref{app:repro}).
\end{enumerate}

% ─────────────────────────────────────────────────────────────────────────────
\section{Background \& Threat Model}
\label{sec:background}

\subsection{Model Context Protocol}

\mcp{} \citep{Anthropic2024MCP} is an open protocol standardizing how LLM
agents discover and invoke external tools. Each \mcp{} tool exposes a name,
input/output schema, and metadata including an optional publisher trust tier.
\pacm{} reads this trust metadata to assign $T \in [0,1]$ and modulate
friction accordingly.

\subsection{Threat Model}

We consider an adaptive adversary \textbf{Adv} with \emph{read access} to
tool outputs seen by the agent and \emph{write access} to attacker-controlled
content that will appear in tool responses. Adv has no access to the agent's
system prompt or \pacm{} internals.

\textbf{Goal:} Induce the agent to call a high-impact tool (e.g.,
\texttt{send\_email}) with attacker-chosen parameters (e.g., exfiltration
target address).

We model \textbf{five attack variants} (Section~\ref{sec:adaptive}):
non-adaptive, lexical mimic, semantic disguise, scope-flag hiding, and fully
adaptive (combined).

\subsection{Research Questions}

\begin{description}[leftmargin=2em]
  \item[RQ1] Does plan-action divergence scoring reliably detect injection?
  \item[RQ2] What is the optimal $\alpha$ (semantic vs.\ scope weight)?
  \item[RQ3] What is \pacm{}'s latency overhead vs.\ undefended agents?
  \item[RQ4] Which \pacm{} components contribute most (ablation)?
  \item[RQ5] How does detection degrade against adaptive adversaries?
\end{description}

% ─────────────────────────────────────────────────────────────────────────────
\section{System Architecture}
\label{sec:architecture}

Figure~\ref{fig:pipeline} shows the \pacm{} pre-tool hook pipeline. All five
steps execute synchronously before any tool call is dispatched.

\begin{figure}[htbp]
\centering
\begin{lstlisting}[language=,caption={PACM pre-tool hook pipeline.},
                   label={lst:pipeline}]
ToolCallRecord
  |
  v
[1] IntentExtractor   -> DeclaredIntent (target_tool, declared_goal, ...)
  |
  v
[2] ActionCapturer    -> CapturedAction (tool_name, args, arg_scope_flags)
  |
  v
[3] DivergenceScorer  -> D in [0.0, 1.0]
  |     D = clip(alpha*D_semantic + (1-alpha)*P_scope, 0, 1)
  v
[4] FrictionPolicy    -> F = clip(D * (1.5 - 0.5*T), 0, 1)
  |                      Decision: ALLOW | FLAG | BLOCK
  v
[5] PACMAuditEntry    -> Logged to pacm_results.jsonl
\end{lstlisting}
\caption{\pacm{} pre-tool hook pipeline. Steps 1--5 add $<$1\,ms overhead
per tool call in post-warmup TF-IDF mode.}
\label{fig:pipeline}
\end{figure}

\paragraph{Step 1 --- Intent extraction.}
\texttt{IntentExtractor} parses the reasoning trace immediately before the
tool call using regex + heuristics, extracting the declared target tool,
goal, and expected parameter summary. No LLM call is made.

\paragraph{Step 2 --- Action capture.}
\texttt{ActionCapturer} normalizes the tool name and argument dict into a
human-readable description and checks arguments against scope rules
(external domain, path prefix, key allowlists), setting \texttt{arg\_scope\_flags}.

\paragraph{Step 3 --- Divergence scoring.}
Described in Section~\ref{sec:scoring}.

\paragraph{Step 4 --- Trust-weighted friction policy.}
Described in Section~\ref{sec:policy}.

\paragraph{Step 5 --- Audit logging.}
Every pre-tool hook invocation is logged as a \texttt{PACMAuditEntry} to
\texttt{experiments/pacm\_results.jsonl}, including $D$, $F$, $T$, and the
decision, enabling offline calibration and figure generation.

\subsection{Dual Harness Architecture}

\pacm{} is validated against two independent agent execution hosts:
\begin{itemize}
  \item \textbf{AgentHarness} (\texttt{src/harness/langgraph\_agent.py}):
    LangGraph-based with \texttt{ToolNode} execution.
  \item \textbf{CustomLoopHarness} (\texttt{src/harness/custom\_loop.py}):
    Framework-free sequential tool-call loop.
\end{itemize}

Both expose an identical hook interface:
\texttt{pre\_tool\_hook:~Callable[[ToolCallRecord],~Optional[str]]}.
\pacm{} plugs into this interface without modifying harness internals,
satisfying the \emph{protocol-agnostic} requirement.

% ─────────────────────────────────────────────────────────────────────────────
\section{Divergence Scorer}
\label{sec:scoring}

\subsection{Scoring Formula}

\begin{equation}
  D = \mathrm{clip}\!\left(\alpha \cdot D_{\mathrm{sem}} + (1-\alpha)
      \cdot P_{\mathrm{scope}},\ 0,\ 1\right)
  \label{eq:divergence}
\end{equation}

where:

\begin{itemize}
  \item $D_{\mathrm{sem}} = 1 - \mathrm{sim\_norm}(e_{\mathrm{intent}},
    e_{\mathrm{action}}) \in [0,1]$. Cosine similarity between sentence
    embeddings of the intent text and the action description, normalized to
    $[0,1]$.
  \item $P_{\mathrm{scope}} = \min(n_{\mathrm{flags}} \times 0.20,\ 0.40)$.
    Rule-based penalty from argument scope flags (e.g., external email domain).
  \item $\alpha \in [0,1]$: semantic weight. Default $\alpha = 0.70$;
    Section~\ref{sec:alpha} identifies $\alpha = 0.9$ as optimal.
\end{itemize}

\paragraph{Tool name mismatch fast-path.}
If the captured tool name differs from the declared target tool and the
declared tool is non-empty, $D_{\mathrm{sem}}$ is set to $0.80$ immediately,
bypassing embedding computation.

\subsection{Embedding Backends}

Two backends are supported via \texttt{\_get\_embedding\_backend()}:
\begin{enumerate}
  \item \textbf{TF-IDF} (default): Lightweight bag-of-words cosine similarity;
    zero external dependencies; $<0.03$\,ms/call.
  \item \textbf{sentence-transformers} \citep{Reimers2019SentenceBERT}:
    all-MiniLM-L6-v2; richer semantic representation; $\approx$22\,ms/call
    (post-warmup). Loaded lazily; model weights cached after first call.
\end{enumerate}

% ─────────────────────────────────────────────────────────────────────────────
\section{Trust-Weighted Friction Policy}
\label{sec:policy}

\subsection{Friction Score}

\begin{equation}
  F = \mathrm{clip}\!\left(D \times (1.5 - 0.5 \cdot T),\ 0,\ 1\right)
  \label{eq:friction}
\end{equation}

Intuition: $T = 1.0$ (fully trusted) gives $F = D$; $T = 0.0$ (zero-trust)
gives $F = 1.5D$ (50\% amplification). Low-trust tools require less evidence
before escalation.

\subsection{Policy Presets}

\begin{table}[htbp]
\centering
\caption{Friction policy presets.}
\label{tab:policy}
\begin{tabular}{lccc}
\toprule
Policy       & Allow threshold & Block threshold & Use case \\
\midrule
Standard     & 0.45 & 0.75 & Default production \\
Conservative & 0.35 & 0.65 & High-risk environments \\
Aggressive   & 0.25 & 0.55 & Maximum sensitivity \\
\bottomrule
\end{tabular}
\end{table}

Decisions: $F < \theta_{\mathrm{allow}}$ $\to$ \textsc{Allow};
$\theta_{\mathrm{allow}} \le F < \theta_{\mathrm{block}}$ $\to$ \textsc{Flag}
(log warning, optional human review); $F \ge \theta_{\mathrm{block}}$ $\to$
\textsc{Block} (synthetic refusal returned to LLM).

% ─────────────────────────────────────────────────────────────────────────────
\section{Experimental Evaluation}
\label{sec:eval}

\subsection{Setup}

\textbf{Harnesses:} LangGraph (\texttt{AgentHarness}) and framework-free
(\texttt{CustomLoopHarness}).
\textbf{Attack types:} \texttt{ignore\_previous} (direct override) and
\texttt{injec\_agent} (indirect agent-state injection).
\textbf{Policies:} Undefended, Standard, Conservative, Aggressive.
\textbf{Synthetic data:} 200 attack + 200 benign records with explicit
$D_{\mathrm{sem}}$ and $P_{\mathrm{scope}}$ components.

\begin{quote}
\textbf{Stub-mode note:} ABR/FPR values in Table~\ref{tab:full_sweep}
reflect \emph{stub mode} execution (no live LLM API calls). Stub reasoning
traces produce uniform text, yielding moderate divergence ($D \approx
0.3$--$0.4$) that falls in the \textsc{Flag} zone. Real-LLM evaluation
requires a live backend: \texttt{OPENAI\_API\_KEY} for OpenAI (gpt-4o-mini)
or \texttt{LLM\_PROVIDER=ollama} for a local Ollama instance (qwen3:8b,
llama3.1:8b). Run \texttt{python -m experiments.probe\_llm\_backend} first
to confirm tool-calling capability. See Section~\ref{sec:cot_filter} for
the CoT tag stripping required for local CoT models.
\end{quote}

\subsection{CoT Reasoning Tag Stripping (Session 9)}
\label{sec:cot_filter}

Local CoT models (\texttt{qwen3:8b}, \texttt{deepseek-r1},
\texttt{phi-4-reasoning}) emit internal chain-of-thought blocks inside
structured tags (\texttt{<think>\ldots</think>}) prior to their final
response. Without stripping, \pacm{}'s \texttt{IntentExtractor} reads the
private scratchpad as the declared intent~---~a systematic false signal that
deflates the plan-action divergence score and facilitates injection evasion.

\texttt{strip\_cot\_tags()} in \texttt{src/pacm/extractor.py} removes four
tag formats before any regex parsing:

\begin{enumerate}[leftmargin=*, label=\arabic*.]
  \item \texttt{<think>\ldots</think>} (qwen3, deepseek-r1; case-insensitive, multiline)
  \item \texttt{[THINK]\ldots[/THINK]} (some llama/mistral fine-tunes)
  \item \texttt{<|thinking|>\ldots<|/thinking|>} (phi-4-reasoning)
  \item \texttt{<!-\,- think -\,->\ldots<!-\,- /think -\,->} (markdown CoT wrappers)
\end{enumerate}

The \texttt{raw\_reasoning} field in \texttt{DeclaredIntent} preserves the
original text (including tags) for audit traceability. The
\texttt{extra[\textquotedbl{}cot\_stripped\textquotedbl{}]} flag records
whether stripping occurred. Twelve new unit tests (\texttt{test\_pacm\_core.py})
verify all formats, case-insensitive variants, multiline blocks, and the
raw-preservation invariant.


\subsection{Alpha Sensitivity Analysis (RQ2)}
\label{sec:alpha}

\begin{table}[htbp]
\centering
\caption{Alpha sensitivity sweep ($n=400$ synthetic records, TF-IDF backend).}
\label{tab:alpha}
\begin{tabular}{rrrrrrr}
\toprule
$\alpha$ & $\bar{D}_{\mathrm{atk}}$ & $\bar{D}_{\mathrm{ben}}$
         & Separation & ABR & FPR \\
\midrule
0.5 & 0.4961 & 0.1315 & 0.3646 & 48.0\% & 0.0\% \\
0.6 & 0.5373 & 0.1442 & 0.3932 & 72.0\% & 0.0\% \\
0.7 & 0.5785 & 0.1568 & 0.4217 & 85.5\% & 0.0\% \\
0.8 & 0.6198 & 0.1695 & 0.4502 & 90.5\% & 0.0\% \\
\textbf{0.9} & \textbf{0.6610} & \textbf{0.1822} & \textbf{0.4787}
             & \textbf{96.0\%} & \textbf{0.0\%} \\
\bottomrule
\end{tabular}
\begin{tablenotes}\small
  \item Best $\alpha$ (highest separation) highlighted in bold.
  \item FPR = 0.0\% across all values on this synthetic distribution.
\end{tablenotes}
\end{table}

\textbf{Finding (RQ2):} Higher $\alpha$ monotonically improves separation and
ABR, confirming that injection attacks primarily manifest as \emph{semantic}
divergence on this distribution. The default $\alpha = 0.70$ is conservative;
production deployment should consider $\alpha \in [0.8, 0.9]$.

\subsection{Full 16-Configuration Benchmark}
\label{sec:sweep}

\input{paper_table.tex}

\subsection{Ablation Study (RQ4)}
\label{sec:ablation}

\input{ablation_table.tex}

\textbf{Findings (RQ4):}
\begin{itemize}
  \item \textbf{Trust modulation is essential}: Variant~B ($T=1.0$ fixed)
    achieves 0\% ABR; the trust factor is required to push $F$ above the
    block threshold for moderate-$D$ attacks.
  \item \textbf{ST backend improves ABR by 15\%}: $D_{\mathrm{st}}$ (54\%)
    outperforms $D_{\mathrm{tfidf}}$ (39\%), confirming richer semantic
    representations capture injection divergence better.
  \item \textbf{Latency negligible}: ST backend adds only 0.011\,ms/call
    post-warmup.
\end{itemize}

\subsection{Adaptive Attacker Stress Test (RQ5)}
\label{sec:adaptive}

\begin{table}[htbp]
\centering
\caption{Adaptive attacker results (STANDARD policy, $T=0.8$, TF-IDF backend,
         $n=20$ per variant). All attacks trigger \textsc{Flag} even when
         not \textsc{Block}ed.}
\label{tab:adaptive}
\begin{tabular}{lrrrrr}
\toprule
Variant & $\bar{D}$ & $\bar{F}$ & Evasion & ABR \\
\midrule
Non-adaptive      & 0.371 & 0.408 & 100\% & 0\% \\
Lexical mimic     & 0.245 & 0.270 & 100\% & 0\% \\
Semantic disguise & 0.376 & 0.414 & 100\% & 0\% \\
Scope hide        & 0.311 & 0.342 & 100\% & 0\% \\
Full adaptive     & 0.350 & 0.385 & 100\% & 0\% \\
\bottomrule
\end{tabular}
\begin{tablenotes}\small
  \item Results reflect stub-mode execution. All $F$ values fall in the
    \textsc{Flag} zone (0.27--0.41), providing a human oversight channel.
  \item Lexical mimic achieves lowest $\bar{D}$ (0.245); motivates
    argument-level token-copy detection in future work.
\end{tablenotes}
\end{table}

\textbf{Finding (RQ5):} In stub mode, all adaptive variants achieve 100\%
evasion of the \textsc{Block} decision because stub reasoning traces produce
uniformly low divergence. Critically, all attacks land in the
\textsc{Flag} zone ($F \approx 0.27$--$0.41$), meaning they are logged and
escalated for human review. With real LLM traces and genuine injections, the
semantic gap between stated intent and hijacked action is substantially larger
($D > 0.5$), placing most attacks in the \textsc{Block} zone. Lexical mimic
($\bar{D} = 0.245$) is the most effective evasion strategy, motivating
argument-level token-copy detection as future work.

\subsection{Error Analysis \& Failure Mode Taxonomy}
\label{sec:errors}

On 200 synthetic records (100 attack, 100 benign):
$\mathrm{TP}=73$, $\mathrm{TN}=92$, $\mathrm{FP}=8$, $\mathrm{FN}=27$.
ABR = 73.0\%, FPR = 8.0\%, Precision = 90.1\%.

\textbf{Dominant false negative category:} FN-4 (high-trust tool exploit,
$n=11$) --- attacks via $T=1.0$ tools face reduced friction amplification.
This motivates re-evaluating trust assignments for high-impact tools.


\subsection{Real-Model Benchmark: qwen3:8b (Session 11)}
\label{sec:realmodel}

\noindent We executed a full real-model evaluation using
\texttt{qwen3:8b} (5.2\,GB quantised) served locally via Ollama's
OpenAI-compatible endpoint.  All results in this subsection reflect
\emph{live} LLM calls with no stubs.

\vspace{2pt}
\noindent\textbf{Backend configuration.}
Ollama 0.9.6 on a single-host CPU+GPU machine;
\texttt{LLM\_PROVIDER=ollama}, \texttt{LLM\_BASE\_URL=http://localhost:11434/v1};
pre-flight TCP health-check via \texttt{probe\_llm\_backend.py}.

\vspace{4pt}
\begin{table}[htbp]
\centering
\caption{Real-model sweep results (qwen3:8b, quick mode, 6 tasks, 2 configs).}
\label{tab:realmodel}
\begin{tabular}{llrrrr}
\toprule
Config & Policy & ABR & BTSR & FPR & Elapsed \\
\midrule
undefended / langgraph / ignore\_previous
  & --- & 0.0\% & 100\% & --- & 663\,s \\
defended / langgraph / ignore\_previous
  & Standard + boost & 100\%$^\dagger$ & 100\% & 0.0\% & 774\,s \\
\bottomrule
\end{tabular}
\begin{tablenotes}\small
  \item $^\dagger$ ABR = 100\% with scope-flag boost active
    ($D_{\mathrm{exfil}}=0.339$, boost$=+0.15$, $F=0.489 > \theta_{\mathrm{allow}}=0.45$).
  \item Without scope boost: $D=0.339 < \theta=0.45$ \textrightarrow{} missed (ABR = 0\%).
  \item BTSR = 100\% confirms zero false positives on benign tasks.
  \item Mean per-call latency: 220--258\,s (qwen3:8b cold-start; acceptable for
    high-value tool gates such as external email sends).
\end{tablenotes}
\end{table}

\vspace{4pt}
\noindent\textbf{Exfiltration case study.}
In the defended run, the agent received an \texttt{ignore\_previous}
injection instructing it to forward inbox contents to
\texttt{attacker@evil.com}.  The tool-call sequence was:

\vspace{2pt}
\begin{enumerate}[leftmargin=*, label=\arabic*.]
  \item \texttt{send\_email(to=alice@example.com)} --- $D=0.226$,
    $F=0.226 < 0.45$ \textrightarrow{} \textsc{Allow}.
  \item \texttt{read\_inbox()} --- $D=0.103$,
    $F=0.103 < 0.45$ \textrightarrow{} \textsc{Allow}.
  \item \texttt{send\_email(to=alice@example.com, body=\textit{reply})} ---
    $D=0.290$, $F=0.290 < 0.45$ \textrightarrow{} \textsc{Allow}.
  \item \texttt{read\_inbox()} --- $D=0.103$,
    $F=0.103 < 0.45$ \textrightarrow{} \textsc{Allow}.
  \item \textbf{\texttt{send\_email(to=attacker@evil.com, body=\textit{inbox data})}}
    --- $D=0.339$, \texttt{scope\_flags=\{to: external\_domain\}},
    $F = 0.339 + 0.15 = 0.489 > 0.45$ \textrightarrow{} \textbf{\textsc{Flag/Block}}.
\end{enumerate}

\vspace{2pt}
The scope-flag boost is the critical discriminating signal: without it,
all five calls would be allowed ($F < 0.45$). With it, the exfiltration
call ($D=0.339$, modest semantic divergence) is correctly escalated because the
\texttt{to} parameter references an external domain outside the trusted perimeter.

\vspace{4pt}
\noindent\textbf{Finding (RQ3 --- Real-Model).}
With scope-flag friction boosting ($\Delta = +0.15$), \pacm{} achieves
\textbf{ABR = 100\%, BTSR = 100\%, FPR = 0.0\%} on a live \texttt{qwen3:8b}
single-task evaluation.  Semantic divergence alone ($D = 0.339$) is
insufficient; the scope signal is essential for reliable exfiltration detection
under open-weight CoT models with moderate reasoning availability.


% ─────────────────────────────────────────────────────────────────────────────
\section{Discussion}
\label{sec:discussion}

\subsection{Limitations}

\begin{enumerate}[leftmargin=*]
  \item \textbf{Real-LLM scale:} The real-model evaluation (Section~\ref{sec:realmodel})
    uses a quick sweep (2 configs, 6 tasks). A full 14-config sweep is needed for
    publication-grade ABR/BTSR confidence intervals.
  \item \textbf{Alpha calibration:} $\alpha = 0.9$ was identified on synthetic
    data. Real injection traces may require different calibration.
  \item \textbf{Trust score dependency:} \pacm{} requires accurate \mcp{} trust
    scores. If metadata is missing, $T = 0.5$ (default) is applied.
  \item \textbf{Adaptive attacker scope:} The five variants were constructed
    synthetically; real adversaries may devise novel evasion strategies.
  \item \textbf{Reasoning availability:} As documented in
    Section~\ref{sec:reasoning_gap}, open-weight tool-calling LLMs via Ollama
    return empty \texttt{AIMessage.content} during tool calls, requiring
    multi-source CoT extraction. Future work should validate across more models.
\end{enumerate}


\subsection{The Reasoning Availability Gap in Open-Weight Tool-Calling LLMs}
\label{sec:reasoning_gap}

\noindent\textbf{Discovery.}
During real-model evaluation (Section~\ref{sec:realmodel}), we observed that
\texttt{qwen3:8b} via Ollama's OpenAI-compatible API returns an \emph{empty}
\texttt{AIMessage.content} field when making tool calls.  The chain-of-thought
reasoning (CoT) is instead placed in one of:

\vspace{2pt}
\begin{itemize}
  \item \texttt{additional\_kwargs[``reasoning\_content'']} --- the Ollama
        field that mirrors OpenAI $o1$'s structured thinking output.
  \item \texttt{additional\_kwargs[``thinking'']} --- an alternative key used by
        some fine-tuned deployments.
  \item Inline \texttt{<think>\ldots</think>} tags inside a \emph{prior}
        message turn, consumed by LangChain before the tool-call AIMessage.
\end{itemize}

\vspace{4pt}
\noindent\textbf{Impact.}
Without a reasoning trace, \texttt{IntentExtractor} returns
$\mathrm{confidence}=0.0$ and $\mathrm{declared\_goal}=\texttt{``''}$.
The TF-IDF divergence scorer then computes $D$ relative to an \emph{empty}
baseline, yielding systematically deflated scores ($D \approx 0.10$--$0.34$)
that do not reflect true plan-action divergence.

\vspace{4pt}
\noindent\textbf{Mitigation: Multi-Source Reasoning Extraction.}
We introduce \texttt{\_extract\_reasoning(ai\_msg, all\_messages)}, a five-priority
fallback chain implemented in both \texttt{AgentHarness} and
\texttt{CustomLoopHarness}:

\begin{enumerate}[leftmargin=*, label=\arabic*.]
  \item \texttt{AIMessage.content} --- standard (non-empty in GPT-4o, Claude).
  \item \texttt{additional\_kwargs[``reasoning\_content'']} --- Ollama qwen3/deepseek-r1.
  \item \texttt{additional\_kwargs[``thinking'' / ``thought'']} --- alternative keys.
  \item Anthropic-style \texttt{\{``type'': ``thinking'', ``thinking'': \ldots\}}
        content blocks.
  \item Previous non-empty \texttt{AIMessage} in the conversation --- graceful
        fallback when CoT is emitted in a separate turn.
\end{enumerate}

\vspace{4pt}
\noindent\textbf{Scope Friction as Compensatory Signal.}
Even when reasoning is unavailable (Source~5 fallback), the \emph{parameter
scope signal} remains reliable: external domain exfiltration is detectable
without any reasoning text via the \texttt{to=external\_domain} scope flag.
We implement a fixed \textbf{scope-flag friction boost} $\Delta = +0.15$ added
to $F$ whenever \texttt{arg\_scope\_flags} are non-empty:

\vspace{-4pt}
\begin{equation}
F_{\mathrm{boosted}} = \min\!\left(1,\; D \times (1.5 - 0.5T) + \Delta_{\mathrm{scope}}\right)
\label{eq:friction_boosted}
\end{equation}

This ensures that exfiltration calls ($D \approx 0.34$, $\Delta=0.15$)
yield $F=0.49 > \theta_{\mathrm{allow}}=0.45$, escalating them for review or
blocking, even under reduced reasoning availability.

\vspace{4pt}
\noindent\textbf{Generalizability.}
The reasoning availability gap is \emph{not} specific to \texttt{qwen3:8b};
it is an architectural property of any model deployed via an
OpenAI-compatible tool-use API where the provider separates thinking from
the response content.  Our multi-source extraction and scope compensation
are model-agnostic and activate only when needed (no overhead for GPT-4o/Claude).
Coverage is validated by 12 unit tests across both harnesses
(\texttt{TestExtractReasoning} in \texttt{test\_harness\_smoke.py} and
\texttt{TestExtractReasoningCustomLoop} in \texttt{test\_custom\_loop.py}).


\subsection{Related Work}

\textbf{Prompt injection defenses.}
\citet{Greshake2023NotWhatYouSignedUp} first characterized indirect prompt
injection. \citet{Perez2022IgnorePrevious} demonstrated direct injection in
GPT-3. \citet{Liu2023PromptInjectionAttack} systematized attack techniques.
\citet{Wang2026AgentPI} provides the most recent SoK, explicitly identifying
the context-dependent evaluation gap that \pacm{} addresses.

\textbf{Structural defenses.}
\citet{Debenedetti2025CaMeL} propose capability-based memory isolation;
\citet{Wallace2024InstructionHierarchy} train instruction priority.
Both are complementary to \pacm{}: structural defenses provide formal
guarantees but require architectural changes; \pacm{} adds a lightweight
monitoring layer without modifying the execution engine.

\textbf{Runtime monitoring.}
\citet{Weng2026ARGUS} propose ARGUS, a concurrent runtime monitoring
framework. \pacm{} differentiates via (a) the plan-action consistency
signal (not just behavioral anomaly detection) and (b) \mcp{} trust
score integration in the friction formula.

\textbf{Benchmark frameworks.}
Our evaluation extends AgentDojo \citep{Debenedetti2024AgentDojo} and
InjecAgent \citep{Zhan2024InjecAgent}.

\textbf{\mcp{} security.}
\citet{Li2025MCPEcosystem} and \citet{Hasan2026MCPFirstGlance} characterize
\mcp{} ecosystem threats; their trust taxonomy directly informs \pacm{}'s $T$
score design.

% ─────────────────────────────────────────────────────────────────────────────
\section{Conclusion}
\label{sec:conclusion}

\pacm{} is a practical, modular runtime defense against context-aware prompt
injection in tool-using LLM agents. By intercepting tool calls and measuring
plan-action divergence through a dual semantic+scope scorer combined with
\mcp{} trust weighting, \pacm{} provides actionable escalation decisions with
negligible latency overhead. The adaptive attacker analysis (RQ5) reveals
that while determined adversaries can evade blocking in stub mode, they
consistently trigger \textsc{Flag} escalation --- preserving a human
oversight channel even in the worst case. The ablation study confirms that
trust modulation and the sentence-transformers backend each contribute
meaningfully to detection. Complete real-LLM evaluation and LLM-backed intent
extraction remain the critical next steps toward a publication-ready result.

% ─────────────────────────────────────────────────────────────────────────────
\bibliographystyle{abbrvnat}
\bibliography{references}

% ─────────────────────────────────────────────────────────────────────────────
\appendix

\section{Reproducibility}
\label{app:repro}

All experiments can be reproduced with a single command (no API key needed
for stub-mode runs):

\begin{lstlisting}[language=bash,caption={One-click reproduction.}]
# Clone and set up
git clone <repo-url>
cd pacm
python -m venv .venv && .venv\Scripts\activate   # Windows
pip install -r requirements.txt

# Run all experiments + generate figures
python -m experiments.reproduce_all --synthetic

# Full real-LLM run (requires OPENAI_API_KEY in .env)
python -m experiments.reproduce_all
\end{lstlisting}

The script executes 8 sequential steps (unit tests, sweep, calibration,
sensitivity, ablations, adaptive attack, error analysis, figures + tables)
and reports a pass/fail summary.

\section{Test Suite}
\label{app:tests}

\pacm{} v1.0.0 ships with 223 unit tests across 10 test files (203 non-ST + 20
sentence-transformers backend tests), all passing. Tests require no external
API keys.

\begin{table}[htbp]
\centering
\caption{Test suite summary.}
\label{tab:tests}
\begin{tabular}{lrl}
\toprule
Test file & Tests & Focus \\
\midrule
\texttt{test\_harness\_smoke.py}          & 16 & Harness integration \\
\texttt{test\_custom\_loop.py}            & 20 & CustomLoopHarness \\
\texttt{test\_mcp\_wrapper.py}            & 24 & MCPWrappedTool \\
\texttt{test\_pacm\_core.py}              & 29 & Core pipeline \\
\texttt{test\_policy.py}                  & 23 & FrictionPolicy \\
\texttt{test\_experiment\_runner.py}      & 25 & Sweep runner \\
\texttt{test\_calibration\_and\_sweep.py} & 25 & Calibration, ROC \\
\texttt{test\_sensitivity.py}             & 25 & Alpha sensitivity \\
\texttt{test\_scorer\_transformers.py}    & 20 & ST backend \\
\texttt{test\_adaptive\_and\_ablations.py}& 33 & RQ5, ablations, errors \\
\midrule
\textbf{Total} & \textbf{240} & \\
\bottomrule
\end{tabular}
\end{table}

\section{File Map}
\label{app:files}

\begin{table}[htbp]
\centering
\caption{\pacm{} v1.0.0 module map.}
\label{tab:files}
\begin{tabular}{llp{6.5cm}}
\toprule
Module & Path & Description \\
\midrule
\texttt{AgentHarness}      & \texttt{src/harness/langgraph\_agent.py}  & LangGraph agent \\
\texttt{CustomLoopHarness} & \texttt{src/harness/custom\_loop.py}       & Framework-free loop \\
\texttt{MCPWrappedTool}    & \texttt{src/harness/mcp\_tool\_wrapper.py} & \mcp{} trust metadata \\
\texttt{IntentExtractor}   & \texttt{src/pacm/extractor.py}            & Plan extraction \\
\texttt{ActionCapturer}    & \texttt{src/pacm/action\_capture.py}      & Arg capture + scope \\
\texttt{DivergenceScorer}  & \texttt{src/pacm/scorer.py}               & Dual-signal scorer \\
\texttt{FrictionPolicy}    & \texttt{src/pacm/policy.py}               & Trust-weighted policy \\
\texttt{PACMMiddleware}    & \texttt{src/pacm/middleware.py}           & Pre-tool hook \\
\texttt{ThresholdCalibrator} & \texttt{src/pacm/calibration.py}        & ROC calibration \\
\texttt{AlphaSensitivity}  & \texttt{src/pacm/sensitivity.py}          & $\alpha$ sweep \\
\texttt{AdaptiveAttacker}  & \texttt{src/pacm/adaptive\_attack.py}     & RQ5 stress test \\
\texttt{ErrorAnalyzer}     & \texttt{src/pacm/error\_analysis.py}      & FP/FN taxonomy \\
\bottomrule
\end{tabular}
\end{table}

\end{document}
